\subsection{FACETS}

The set of points of E that do not belong to any hyperplane H of the set H is open since H is locally finite. More precisely, we have the following result:

PROPOSITION 1. Let a be a point of E. There exists a connected open neighbourhood of a that does not meet any hyperplane H that belongs to H and does not pass through a. Moreover, there exist only finitely many hyperplanes that belong to H and pass through a.

The set N of hyperplanes H such that \(H \in H\) and \(a \notin H\) is locally finite since it is contained in H. Hence, the set U of points of E that do not belong to any of the hyperplanes of the set N is open. Since \(a \in U\) , there is a connected open neighbourhood of a contained in U. The remainder of the proposition is clear.

Given two points \(x\) and \(y\) of \(\mathbf{E}\) , denote by \(\mathrm{R}\{x,y\}\) the relation

``For any hyperplane \(H \in H\) , either \(x \in H\) and \(y \in H\) or x and y are strictly on the same side of H.''

Clearly, R is an equivalence relation on E.

DEFINITION 1. A facet of E relative to \(\mathfrak{H}\) is an equivalence class of the equivalence relation defined above.

PROPOSITION 2. The set of facets is locally finite.

This is clear since \(\mathfrak{H}\) is locally finite.

Let F be a facet and \(a\) a point of F. A hyperplane \(H \in H\) contains F if and only if \(a \in H\) ; the set F of these hyperplanes is thus finite; their intersection is an affine subspace L of E, which we shall call the affine support of F; the dimension of L will be called the dimension of F.

If \(\mathfrak{N}\) is the set of hyperplanes \(H\in \mathfrak{H}\) not containing \(F\) , then

\[
F = L \cap \bigcap_ {H \in \mathfrak {N}} D _ {H} (a).\tag{2}
\]

We shall prove that the closure of F is given by

\[
\overline {{F}} = L \cap \bigcap_ {H \in \mathfrak {N}} \overline {{D _ {H} (a)}}.\tag{3}
\]

It is clear that the right-hand side contains the left-hand side. Conversely, let \(x \in L \cap \bigcap_{H \in \mathfrak{N}} \overline{\mathrm{D}_H(a)}\) . The open segment with extremities \(a\) and \(x\) is contained in \(L\) and in each of the \(D_H(a)\) for \(H \in \mathfrak{N}\) , and hence in \(F\) . It follows that \(x\) is in the closure of \(F\) , hence the formula.

PROPOSITION 3. Let F be a facet and L its affine support.

\begin{enumerate}
\def\labelenumi{(\roman{enumi})}
\item
  The set \(\mathbf{F}\) is a convex open subset of the affine subspace \(\mathbf{L}\) of \(\mathbf{E}\) .
\item
  The closure of \(\mathbf{F}\) is the union of \(\mathbf{F}\) and facets of dimension strictly smaller than that of \(\mathbf{F}\) .
\item
  In the topological space L, the set F is the interior of its closure.
\end{enumerate}

Since every open half-space and every hyperplane are convex subsets of E, formula (2) shows that F is the intersection of a family of convex subsets, and hence is convex. On the other hand, let a be in F, and let U be a convex open neighbourhood of a in E that does not meet any of the hyperplanes in the set N of \(H \in H\) such that \(a \notin H\) . For any \(H \in N\) , we thus have \(U \subset D_{H}(a)\) , hence \(L \cap U \subset F\) , so that F is open in the topological space L.

Let \(b\) be a point of \(\overline{F} - F\) , belonging to a facet \(F'\) , and let \(\mathfrak{N}'\) be the set of hyperplanes \(H \in \mathfrak{N}\) passing through \(b\) ; put \(\mathfrak{N}'' = \mathfrak{N} - \mathfrak{N}'\) . For any \(H\) in \(\mathfrak{N}''\) , we have \(b \notin H\) and \(b \in \overline{D_H(a)}\) , hence \(b \in D_H(a)\) and \(D_H(b) = D_H(a)\) ; by the definition of a facet, we thus have

\[
F ^ {\prime} = L \cap \bigcap_ {H \in \mathfrak {N} ^ {\prime}} H \cap \bigcap_ {H \in \mathfrak {N} ^ {\prime \prime}} D _ {H} (a),\tag{4}
\]

whereas (3) implies that

\[
\overline {{\mathbf {F}}} = \mathbf {L} \cap \bigcap_ {\mathsf {H} \in \mathfrak {N} ^ {\prime}} \overline {{\mathsf {D} _ {\mathsf {H}} (a)}} \cap \bigcap_ {\mathsf {H} \in \mathfrak {N} ^ {\prime \prime}} \overline {{\mathsf {D} _ {\mathsf {H}} (a)}},\tag{5}
\]

hence \(\mathbf{F}' \subset \overline{\mathbf{F}}\) . We cannot have \(\mathfrak{N}' = \varnothing\) , for this would imply that \(\mathbf{F} = \mathbf{F}'\) by (2) and (4), contrary to the hypothesis that \(b \notin \mathbf{F}\) and \(b \in \mathbf{F}'\) . The support of \(\mathbf{F}'\) is the set \(\mathbf{L}' = \mathbf{L} \cap \bigcap_{\mathbf{H} \in \mathfrak{N}'} \mathbf{H}\) ; we have \(a \in \mathbf{L}\) , but \(a \notin \mathbf{H}\) for \(\mathbf{H}\) in \(\mathfrak{N}'\) , so \(\mathbf{L}' \neq \mathbf{L}\) and finally \(\dim \mathbf{L}' < \dim \mathbf{L}\) , that is \(\dim \mathbf{F}' < \dim \mathbf{F}\) . This proves (ii).

Let H be in \(N'\) and let D be the open half-space bounded by H and distinct from \(D_{\mathrm{H}}(a)\) ; we have \(b \in H \cap L\) , and it is immediate that \(D \cap L\) is a half-space of L bounded by the hyperplane \(H \cap L\) of L; consequently, every neighbourhood of b in L meets \(D \cap L\) , and since \(D \cap L\) is disjoint from \(\overline{F}\) by (3), we see that the point b of \(\overline{F} - F\) cannot be in the interior of \(\overline{F}\) in the topological space L. Since F is open in L, we have (iii). Q.E.D.

COROLLARY. Let F and \(F'\) be two facets. If \(\overline{F} = \overline{F'}\) , the facets F and \(F'\) are equal.

This follows from (iii).

PROPOSITION 4. Let F be a facet, and let L be an affine subspace of E that is the intersection of hyperplanes belonging to H: denote by N the set of hyperplanes \(H \in H\) that do not contain L. The following conditions are equivalent:

\begin{enumerate}
\def\labelenumi{(\roman{enumi})}
\item
  There exists a facet \(\mathbf{F}'\) with support L that meets \(\overline{\mathbf{F}}\) .
\item
  There exists a facet \(\mathbf{F}'\) with support \(\mathbf{L}\) contained in \(\overline{\mathbf{F}}\) .
\item
  There exists a point \(x\) in \(\mathbf{L} \cap \overline{\mathbf{F}}\) that does not belong to any of the hyperplanes of \(\mathfrak{N}\) .
\end{enumerate}

If these conditions are satisfied, \(\mathrm{L} \cap \mathrm{D}_{\mathfrak{M}}(\mathrm{F})\) is the unique facet with support \(\mathrm{L}\) contained in \(\overline{\mathbf{F}}\) .

\begin{enumerate}
\def\labelenumi{(\roman{enumi})}
\item
  \(\Longrightarrow\) (ii): Since \(\overline{\mathbf{F}}\) is a union of facets (Prop. 3, (ii)), every facet that meets \(\overline{\mathbf{F}}\) meets a facet contained in \(\overline{\mathbf{F}}\) , and so is equal to it.
\item
  \(\Longrightarrow\) (iii): Every point \(x\) of \(F'\) satisfies (iii) since every hyperplane of \(\mathfrak{H}\) containing \(x\) contains \(F'\) , and hence \(L\) .
\item
  \(\Longrightarrow\) (i): Let \(x\) be a point satisfying (iii) and let \(\mathbf{F}'\) be the facet containing \(x\) ; it is clear that \(\mathbf{F}'\) meets \(\overline{\mathbf{F}}\) . Let \(\mathrm{H} \in \mathfrak{H}\) ; then \(x \notin \mathrm{H}\) if \(\mathrm{H} \in \mathfrak{N}\) and clearly \(x \in \mathrm{H}\) if \(\mathrm{H} \notin \mathfrak{N}\) ; consequently the support of \(\mathbf{F}'\) is the intersection of the hyperplanes of \(\mathfrak{H} - \mathfrak{N}\) , and is equal to \(\mathrm{L}\) .
\end{enumerate}

Finally, let \(F'\) be a facet with support L contained in \(\overline{F}\) , and let x be a point of \(F'\) ; since no hyperplane of \(N \subset H\) passes through x, Prop. 1 shows that there exists a convex open neighbourhood U of x that does not meet any hyperplane of N. Since x is in the closure of F, \(U \cap F \neq \varnothing\) ; now N is the set of hyperplanes \(H \in H\) that do not contain \(F'\) , and for all H in N we have \(D_{H}(x) = D_{H}(U) = D_{H}(U \cap F) = D_{H}(F)\) and formula (2) implies that

\[
\mathrm{F} ^ {\prime} = \mathrm{L} \cap \mathrm{D} _ {\mathfrak {N}} (\mathrm{F}). \quad \text { Q.E.D. }
\]

